\documentclass[10pt]{article}
\usepackage[T1]{fontenc}
\usepackage[utf8]{inputenc}
\usepackage[french]{babel}
\usepackage{amsmath,amssymb,mathrsfs,mathpazo}
\usepackage{geometry}
\geometry{a4paper,margin=2cm}
\usepackage{graphicx,array,tabularx,tabularray,colortbl,multirow,multicol,booktabs,enumitem,pifont,ulem,fancybox,diagbox,cancel}
\usepackage{pstricks,pstricks-add,pst-plot,pst-tree,pst-node,pst-text,pst-func,pst-math,pst-bspline,pst-3dplot,pst-eucl,pst-solides3d}
\let\pscalculate\undefined
\usepackage{pst-bezier}
\usepackage{tikz,pgfplots,tkz-tab}
\usetikzlibrary{arrows,arrows.meta,patterns,calc,positioning,trees,decorations.markings,shapes.geometric,shapes.misc,shapes.arrows,intersections,decorations.pathmorphing,decorations.pathreplacing,backgrounds,shadows,babel}
\pgfplotsset{compat=1.18}
\usepgfplotslibrary{fillbetween}
\usepackage[np]{numprint}
\usepackage{listings,makecell,pas-tableur,twoopt}
\let\barre\undefined
\usepackage[pstricks]{tablvar}
\definecolor{darkgreen}{RGB}{0,100,0}
\providecommand{\red}{\color{red}}
\providecommand{\blue}{\color{blue}}
\providecommand{\green}{\color{green}}
\providecommand{\black}{\color{black}}
\providecommand{\white}{\color{white}}
\providecommand{\gray}{\color{gray}}
\providecommand{\cyan}{\color{cyan}}
\providecommand{\magenta}{\color{magenta}}
\providecommand{\yellow}{\color{yellow}}
\providecommand{\violet}{\color{violet}}
\providecommand{\orange}{\color{orange}}
\providecommand{\euro}{\texteuro}
\providecommand{\EUR}{\texteuro}
\providecommand{\ieme}{\textsuperscript{e}}
\providecommand{\degres}{\ensuremath{^\circ}}
\providecommand{\degree}{\ensuremath{^\circ}}
\providecommand{\vv}[1]{\overrightarrow{#1}}
\providecommand{\Coord}[1]{#1}
\providecommand{\decosix}{}
\providecommand{\starredbullet}{\ensuremath{\bullet}}
\pagestyle{empty}
\setlength{\parindent}{0pt}
\usepackage{adjustbox}\DeclareUnicodeCharacter{2236}{:}\DeclareUnicodeCharacter{2190}{\ensuremath{\leftarrow}}
\providecommand{\R}{\mathbb{R}}
\providecommand{\N}{\mathbb{N}}
\providecommand{\Z}{\mathbb{Z}}
\providecommand{\Q}{\mathbb{Q}}
\providecommand{\D}{\mathbb{D}}
\providecommand{\C}{\mathbb{C}}
\providecommand{\e}{{\mathrm{e}}}
\providecommand{\vect}[1]{\overrightarrow{#1}}
\providecommand{\vectt}[1]{\overrightarrow{\mathrm{#1}}}
\providecommand{\vvt}[1]{\overrightarrow{\mathrm{#1}}}
\providecommand{\np}[2][]{{#2~#1}}
\providecommand{\vv}[1]{{\overrightarrow{#1}}}
\renewcommand{\d}{{\mathrm{d}}}
\providecommand{\degres}{{^\circ}}
\providecommand{\degre}{{^\circ}}
\providecommand{\blue}{\color{blue}}
\providecommand{\red}{\color{red}}
\providecommand{\enskip}{\quad}
\providecommand{\llbracket}{\lbrack\!\lbrack}
\providecommand{\rrbracket}{\rbrack\!\rbrack}
\providecommand{\euro}{{\text{€}}}
\providecommand{\no}{{\text{n°}}}
\providecommand{\interval}[2]{{[#1;#2]}}
\providecommand{\intervFO}[2]{{[#1;#2[}}
\providecommand{\dotfill}{\ldots}

\providecommand{\ds}{\displaystyle}
\providecommand{\textding}[1]{#1}
\providecommand{\Oij}{\(\mathrm{O};\vec{\imath},\vec{\jmath}\)}
\providecommand{\Oijk}{\(\mathrm{O};\vec{\imath},\vec{\jmath},\vec{k}\)}
\providecommand{\Ouv}{\(\mathrm{O};\vec{u},\vec{v}\)}
\providecommand{\intervFF}[2]{\ensuremath{[#1;#2]}}
\renewcommand{\euro}{\ifmmode\text{\texteuro}\else\texteuro\fi}
\DeclareUnicodeCharacter{00B0}{\ifmmode^{\circ}\else\textdegree\fi}
\DeclareUnicodeCharacter{20AC}{\ifmmode\text{\texteuro}\else\texteuro\fi}
% Énoncés : Éducation nationale. Transcription et corrigés : APMEP, Denis Vergès et contributeurs. Découpage et classement : CoopMaths, Sébastien Lozano, Éric Elter et contributeurs. Adaptation de la présentation : Maths Explorer.
% CC BY-SA : https://coopmaths.fr/www/about/
% https://coopmaths.fr/alea/static/e3c/2021/tex/e3c_2021_01_specimen2_4_cor.tex
\begin{document}
\noindent\textit{APMEP / CoopMaths. CC BY-SA. Présentation adaptée par Maths Explorer.}\par\medskip
\bigskip

Soit $f$ la fonction définie sur $\R$ par:  

\[f(x)=\e^{2x} + 6 \e^{x} - 8x-4.\]

\begin{list}{\textbullet}{Dans le plan rapporté à un repère orthogonal, on considère:}
\item $\mathcal{C}_f$ la courbe représentative de la fonction $f$;
\item $\mathcal{D}$ la droite d'équation cartésienne $y=-8x-4$.
\end{list}



\begin{enumerate}
\item
\begin{list}{\textbullet}{}
\item Pour tout $x \in \R$, $f'(x) = 2\e^{2x} +6\e^{x}-8$.
\item Pour tout $x \in \R$, $2\left (\e^{x}-1\right )\left (\e^{x}+4\right ) = 2\left ( \e^{2x}-\e^{x}+4\e^{x}-4\right ) = 2\e^{2x} +6\e^{x}-8$.
\end{list}

On peut donc en déduire que, pour tout $x \in \R$, $f'(x)=2\left (\e^{x}-1\right )\left (\e^{x}+4\right )$.

\item On étudier le signe de $f'(x)$ sur $\R$.

\begin{list}{\textbullet}{}
\item Pour tout x de $\R$, $\e^{x}>0$ donc $\e^{x}+4>0$.
\item On sait que $\e^{x}>1 \iff x>0$ donc $\e^{x}-1>0$ sur $]0~;~+\infty[$.
\item De plus $\e^{x}-1=0 \iff \e^{x}=1 \iff x=0$.
\end{list}

D'où le tableau de signes de $f'(x)$:

\begin{center}
\renewcommand{\arraystretch}{1.5}
\def\esp{\hspace*{2cm}}
$\begin{array}{|c | *{5}{c} |} 
\hline
x  & -\infty & \esp & 0 & \esp  & +\infty \\
\hline
f'(x) &  & \pmb{-} &  \vline\hspace{-2.7pt}{0} & \pmb{+} &    \\
\hline
\end{array}$
\renewcommand{\arraystretch}{1}
\end{center}

\item La dérivée s'annule et change de signe pour $x=0$; $f(0) = \e^{0} + 6\e^{0} -0 -4 = 3$.

On dresse le tableau de variations de la fonction $f$ sur $\R$.

\begin{center}
{\renewcommand{\arraystretch}{1.3}
\psset{nodesep=3pt,arrowsize=2pt 3}  
\def\esp{\hspace*{1.5cm}}
\def\hauteur{0pt}
$\begin{array}{|c| *4{c} c|}
\hline
 x & -\infty & \esp & 0 & \esp & +\infty \\
 \hline
f'(x) &  &  \pmb{-} & \vline\hspace{-2.7pt}0 & \pmb{+} & \\  
\hline
  & \Rnode{max1}{~}  &  &  &  & \Rnode{max2}{~}   \\
f(x) & &  & & &  \rule{0pt}{\hauteur} \\
 &  & &   \Rnode{min}{3} & & \rule{0pt}{\hauteur}
\ncline{->}{max1}{min} \ncline{->}{min}{max2}\\
\hline
\end{array}$
}
\end{center}

\item 
La fonction $f$ admet pour minimum le nombre strictement positif 3, donc pour tout $x$ de $\R$, $f(x)>0$.

\item La courbe $\mathcal{C}_f$ et la droite $\mathcal{D}$ ont un point commun si et seulement si l'équation $f(x)=-8x-4$ admet une solution.

On résout cette équation:

$f(x)=-8x-4
\iff
\e^{2x} + 6 \e^{x} - 8x-4 = -8x-4
\iff
\e^{2x} + 6 \e^{x} = 0
\iff
\e^{x} \left (\e^{x}+6\right ) = 0$

Pour tout réel $x$, $\e^{x}>0$ donc le produit $\e^{x} \left (\e^{x}+6\right )$ n'est jamais nul.

L'équation n'a donc pas de solution, ce qui veut dire que la courbe $\mathcal{C}_f$ et la droite $\mathcal{D}$ n'ont pas de point commun.



\end{enumerate}


\begin{center}
\large \textbf{Figure (non demandée)}
\end{center}

\bigskip

\begin{center}
\psset{xunit=1cm,yunit=0.2cm,labelFontSize=\scriptstyle}
\def\xmin {-6}   \def\xmax {4}
\def\ymin {-5}   \def\ymax {40}
\begin{pspicture*}(\xmin,\ymin)(\xmax,\ymax)

\psgrid[unit=1cm,subgriddiv=1,  gridlabels=0, gridcolor=lightgray](-6,-1)(4,8)

\psaxes[arrowsize=3pt 2, ticksize=-2pt 2pt,Dy=5]{->}(0,0)(-5.99,-4.99)(\xmax,\ymax) 
\uput[dl](0,0){O}
\psplot[linecolor=blue,plotpoints=5000]{\xmin}{\xmax}{2.7183 2 x mul exp 2.7183 x exp 6 mul add x 8 mul sub 4 sub}
\uput[r](1.43,27.1){\blue $\mathcal{C}_f$}
\psplot[linecolor=red,plotpoints=5000]{\xmin}{\xmax}{-8 x mul 4 sub}
\uput[dl](-1,4){\red $\mathcal{D}$}
\end{pspicture*}
\end{center}

\end{document}
